\documentclass[11pt]{article}
\usepackage[margin=0.85in]{geometry}
\usepackage{amsmath,amssymb,booktabs,longtable,hyperref}
\hypersetup{colorlinks=true,urlcolor=blue}
\setlength{\parindent}{0pt}
\setlength{\parskip}{0.55em}
\title{Executable Observer--Controller Uncertainty Integration}
\author{Collision Avoidance Research}
\date{October 3, 2026}
\begin{document}
\maketitle

\textbf{Result.} The observer's timestamped error enclosures now enter the
controller's prediction and constraints. Constant target-parameter uncertainty
is propagated analytically; ego uncertainty is propagated through the existing
affine model. Collision, path, physical-state, terminal-entry and the single
CLF constraints are tightened. This is a partial implementation of
\path{controller/OBSERVER_ROBUST_PCBF_THEORY.tex}, not completion of its
nonlinear output-feedback safety certificate.

The integration also exposes a concrete blocker: with a valid representative
observer publication, the required terminal uncertainty reserve is about
315 times the existing nominal terminal-core radius. The controller rejects
that problem instead of silently treating the estimated state as exact.

\section{Motion contract and implementation scope}
The true target obeys
\[
 \dot V_C=A_C,\qquad \dot A_C=0,\qquad \dot\beta_C=0,\qquad
 \dot\psi_C=V_C\sin\beta_C/l_r.
\]
There is no target jerk disturbance, sideslip-rate state, or off-contract
maneuver in this study. Updating an estimate of $A_C$ or $\beta_C$ does not
mean that the true parameters change. Every member of a current prediction
set obeys this same motion contract. Existing observer API support for other
research configurations was not extended or used here.

The controller remains one PCBF--CLF controller. Previous inputs normally
initialize the new nonlinear rollout. An unavailable or failed shift can use
the existing Zhai-inspired potential-field initializer. Each anchor is
linearized once; no executable backup or additional CLF was introduced.
The collision node margin remains $0.20\,\mathrm m$, the path-center deviation
limit remains $10\,\mathrm m$, and road boundaries remain disabled.

\section{Observer publication and exact target-set propagation}
\texttt{nrmmControllerErrorBounds} now publishes an aligned relative-position
ball, relative-course bound, and intervals for current speed, constant
acceleration and constant curvature $\kappa_C=\sin\beta_C/l_r$. The inverse
reconstruction's acceleration estimate can be clipped. Its error radius is
therefore recentered at the acceleration actually published to the controller:
\[
 b_{A,\mathrm{published}}
 =b_{A,\mathrm{raw}}+|\widehat A_{\mathrm{raw}}-\widehat A_{\mathrm{published}}|.
\]
The parameter intervals are intersected with the declared physical domain.
Current timestamp, availability, center alignment and finite radii are part
of this interface. Missing or invalid certificates are not converted to zero
uncertainty. Inputs without uncertainty metadata continue to declare exact
states, recovering the existing baseline.

For a current nominal target, write
\[
 s(t)=V_0t+\tfrac12 A_0t^2,\qquad
 b_s(t)=b_V|t|+\tfrac12b_At^2.
\]
The analytic enclosure used by \texttt{targetErrorTube} is
\begin{align*}
 b_p(t)&=b_p(0)+b_s(t)+|s(t)|b_\gamma(0)+\tfrac12s(t)^2b_\kappa,\\
 b_\gamma(t)&=\min\{\pi,b_\gamma(0)+|s(t)|b_\kappa
                       +( |\kappa_0|+b_\kappa)b_s(t)\},\\
 b_\psi(t)&=\min\{\pi,b_\gamma(t)+b_\beta\}.
\end{align*}
To obtain the position bound, split the uncertain arc endpoint from the
integral over the nominal arc. The unit tangent bounds the first difference
by $b_s$; integration of course and curvature differences bounds the second
by $|s|b_\gamma+\tfrac12s^2b_\kappa$. The formulas also enclose the existing
signed-arc convention. No stop clamp was added to alter the target contract.
These bounds do not assume future measurements or shrink a reachable set
using a future observer convergence rate.

\section{Constraint tightening and update semantics}
Pairwise rectangle distance is invariant to a common rigid pose. When the
relative target publication and ego state share the same reference pose,
the collision calculation uses a frame frozen at that sample. Its initial
ego generator retains body-velocity and yaw-rate uncertainty, while common
absolute position and orientation cancel. Global pose uncertainty is still
included in path and CLF constraints.

The ego generator is propagated as $G_{i+1}=A_iG_i$, using the dynamics
Jacobians already computed for the affine optimization. For a fixed
separating normal $n$, the added collision reserve is
\[
 \|n^\top G_{p,i}\|_1+\|n\|_2
 \left[b_{p,C}(t_i)+2R_E\sin\frac{\min(\pi,b_{\psi,E,i})}{2}
                     +2R_C\sin\frac{b_{\psi,C}(t_i)}{2}\right],
\]
where $R_E,R_C$ include rectangle center offsets. The orientation terms are
geometric chord bounds. The ego propagation and nominal collision
convexification remain affine approximations; no whole-domain nonlinear
dynamics or nominal-geometry remainder was added.

Updated target estimates are accepted without requiring the user to reset
the controller. The previous input plan remains a warm start. However,
observer-input frames recompute their primary safety budget: the code has
not established inclusion of a new information set in the previous
successor set. Consequently, an old slack sum cannot certify a cap for the
new set. A fresh current target enclosure is required after uncertain tracking;
silently continuing an old target as exact is not permitted.

Terminal entry reserves
\[
 r_{\mathrm{unc}}=\sum_j\|L_fG_{\mathrm{intrinsic},j}\|_2,
 \qquad \|L_fe_f\|_2\le r_f-r_{\mathrm{unc}}
\]
inside the existing free-pose quotient family. This is a sufficient affine
entry test, not robust invariance of an output-feedback terminal policy.
The nominal terminal separation/departure construction is unchanged.

\section{The same CLF, with bounded uncertainty and a vanishing tie}
The single function remains $V=e^\top Pe$, with desired loss
$\ell=e^\top Qe$ and decrease coefficient $\eta=0.5$. Define
\begin{align*}
 b_0&=\sum_j\|P^{1/2}J_0G_{0,j}\|_2,&
 b_1&=\sum_j\|P^{1/2}J_1G_{1,j}\|_2,\\
 b_Q&=\sum_j\|Q^{1/2}J_0G_{0,j}\|_2,&
 B_0&=(\sqrt{V_0}-b_0)_+^2-\eta(\sqrt{\ell_0}+b_Q)^2.
\end{align*}
For fixed positive $w$, Young's inequality gives the convex sufficient row
\[
 (1+w)V_{1,\mathrm{aff}}+(1+1/w)b_1^2\le B_0+\rho.
\]
With $b_1=0$, the implementation uses the exact nominal quadratic instead.
The existing single SOC represents this inequality. Mean successor values
remain separate from uncertainty upper bounds in model-agreement diagnostics.
The objective is now
\[
 \frac{\rho+\epsilon\underline\ell R(\Delta U)}{\max(1,V_0)},\qquad
 \underline\ell=(\sqrt{\ell_0}-b_Q)_+^2,\quad 0\le R\le1,
 \quad \epsilon=10^{-4}.
\]
The tie penalizes input deviations from the linearization inputs. Its possible
physical CLF-slack sacrifice is at most $\epsilon\underline\ell$ at an exact
optimum and vanishes at nominal behavior. Solver tolerances and nonlinear
approximation errors are additional. This implements the affine sufficient
condition; it does not prove global nonlinear CLF dissipation.

\section{Validation and the measured terminal obstruction}
The complete MATLAB suite passed: \textbf{677 passed, zero failed, zero
incomplete}. Focused tests cover exact target-family enclosure, zero-radius
reduction, common-pose cancellation, changed target estimates, geometric
tightening, affine CLF bounds and unavailable/stale enclosures.

The deterministic interface probe increases the target position radius from
$0.01$ to $0.10\,\mathrm m$. The maximum collision-row reserve increases from
$0.01$ to $0.10\,\mathrm m$, and the future optimized input sequence changes
(Frobenius difference approximately $1.95\times10^{-4}$ in mixed steering/input
coordinates). The first command is effectively unchanged in this particular
probe. A $1\,\mathrm{mm}$ position-estimate correction and a
$0.001\,\mathrm{m/s^2}$ acceleration-estimate correction are accepted using a
shifted input initialization, one linearization and a recomputed primary
budget. Both acceleration intervals contain the same true constant $A=0$.

The real-runtime probe uses a cruising ego at $10\,\mathrm{m/s}$ and a target
at relative position $(20,5)\,\mathrm m$, traveling at $12\,\mathrm{m/s}$ in
the opposite direction with $A=\beta=0$. The points initially equal truth;
declared valid initial error bounds are nonzero. Published ego radii are
approximately $(0.12,0.12,0.002,0.02,0.02,0.002)$ in the controller's state
units. The target parameter intervals are
\[
 V_C\in[11.95,12.05],\quad A_C\in[-0.02,0.02],\quad
 \kappa_C\in[-1.40054,1.40054]\times10^{-4}.
\]
The $4.4\,\mathrm s$ target position radius grows to approximately
$0.949\,\mathrm m$. The interface is accepted, but the robustified affine
optimization has no solution. Its terminal metrics are
\begin{center}
\begin{tabular}{lr}
\toprule
Quantity & Value\\
\midrule
Nominal terminal weighted radius $r_f$ & $0.00048828125$\\
Required affine uncertainty reserve $r_{\mathrm{unc}}$ & $0.1540034$\\
Reserve/radius ratio & $315.399$\\
Pure longitudinal-speed half-width of the core & $0.00006435\,\mathrm{m/s}$\\
Propagated longitudinal-speed radius & $0.0202953\,\mathrm{m/s}$\\
\bottomrule
\end{tabular}
\end{center}
The two weighted radii are not physical distances. Even without a difficult
collision, the fixed-input uncertainty family cannot fit into this tiny
precise-state terminal neighborhood. This is not evidence that the real target
violated its model, nor evidence of physical inevitability. It identifies the
need for a causal prediction policy and an information-state terminal design.

\section{Exact-state scenario regression}
The regression uses the existing fourteen scenarios, speeds $8$ and
$15\,\mathrm{m/s}$, sampling at $0.05\,\mathrm s$, and independent tight
\texttt{ode45} replay. No observer uncertainty is injected in this campaign.
Nominal recovery requires the five existing error tolerances for a one-second
dwell after the prescribed encounter. Frame caps are observation limits,
not success criteria. Thirteen scenarios avoid collision and recover; the
$15\,\mathrm{m/s}$ accelerating head-on case still stops after five executed
holds, at $t=0.25\,\mathrm s$, because model agreement is not reached.

\begin{center}
\begin{tabular}{lrrrr}
\toprule
Scenario & Speed & Minimum gap (m) & Recovery (s) & Outcome\\
\midrule
Head-on & 8 & 1.3921 & 10.20 & Recovered \\
Accelerating head-on & 8 & 1.0645 & 10.25 & Recovered \\
Braking lead & 8 & 0.5901 & 13.15 & Recovered \\
Crossing & 8 & 0.1626 & 12.55 & Recovered \\
Turning crossing & 8 & 0.1936 & 12.80 & Recovered \\
Curved head-on & 8 & 1.7863 & 9.90 & Recovered \\
Curved crossing & 8 & 0.1904 & 9.95 & Recovered \\
Head-on & 15 & 0.2228 & 7.35 & Recovered \\
Accelerating head-on & 15 & 13.4432 & --- & Stopped \\
Braking lead & 15 & 1.2591 & 18.55 & Recovered \\
Crossing & 15 & 1.2889 & 8.60 & Recovered \\
Turning crossing & 15 & 1.2491 & 8.10 & Recovered \\
Curved head-on & 15 & 2.3221 & 63.75 & Recovered \\
Curved crossing & 15 & 1.0060 & 8.00 & Recovered \\
\bottomrule
\end{tabular}
\end{center}
The failed case's minimum gap only describes its executed prefix. The
overall minimum replay gap is approximately $0.16263\,\mathrm m$: positive,
but smaller than the nominal $0.20\,\mathrm m$ node margin. This remains a
demonstration of the gap between affine sampled constraints and nonlinear
whole-hold clearance, not a verified $0.20\,\mathrm m$ robust guarantee.

The extended curved head-on case has a second encounter because both target
and reference continue on circles under their fixed motion laws. The driver's
baseline last-contact time is $56.445\,\mathrm s$. Extending its unfinished
run confirms recovery at $63.75\,\mathrm s$, rather than classifying a time
cap as controller failure. The target contract is not changed to force it to
leave. Timing records are preserved, but overlap with other MATLAB validation
means this campaign is not a dedicated real-time benchmark.

\section{What is still required}
The operational connection is implemented and tested, but these remaining
premises are substantive, not documentation details:
\begin{enumerate}
 \item A causal output-feedback prediction tube must replace the fixed-input
 affine image if future feedback is used to reduce spread. Future estimation
 error must enter with its correct sign and a valid bound.
 \item The terminal information set and its policy must accommodate the
 observer error floor. Enlarging the old radius without a new invariant-set
 argument, or omitting its uncertainty reserve, would not establish this.
 \item Whole-domain nonlinear dynamics/error-chart/geometry remainders and
 whole-hold separation reserves are needed for nonlinear collision guarantees.
 \item Posterior intersection and successor-set inclusion are needed before
 noisy frames can safely inherit a previous PCBF budget.
\end{enumerate}
Positive PCBF slack remains an explicit relaxation, not a requested-clearance
certificate. The current metadata explicitly reports that nonlinear robust
safety and recursive feasibility remain uncertified.

\section{Reproducibility}
Machine-readable tests, integration probes and compact scenario results are
in \path{report/OBSERVER_CONTROLLER_INTEGRATION_20261003/}.
\path{scripts/verifyObserverControllerIntegration.m} regenerates the targeted
interface and terminal-obstruction measurements. The experiment commands,
source hashes and evidence locations are recorded in the associated
\texttt{provenance.json}. Local logs and full scenario continuations are kept
as identified external export copies with the project commit record.
\end{document}
